\section{Goal and Vision}

\subsection{Problem Statement}

Classic \UML{} remains valuable because it is Linux itself compiled
to run as a userspace program. That gives researchers a real kernel,
real system call semantics, and a deployable target in environments
where full hardware virtualization is not available. The problem is
that the historical execution model does not provide the performance,
feature composition, or operational discipline expected by modern
kernel fuzzing and observability workflows.

The redesign sets a higher bar:
\begin{itemize}
  \item production \UML{} should approach native syscall cost when
        the host exposes \KVM{};
  \item research \UML{} should ship with sanitizers, tracing,
        kprobes, ftrace, BPF, KGDB, and debug surfaces composed from
        one tree;
  \item fuzzing \UML{} should provide KCOV, snapshot/forkserver
        startup, and a path toward a syzkaller \texttt{vm/uml}
        backend;
  \item sandbox \UML{} should compile out debug surfaces and minimize
        host attack surface;
  \item time-travel and record/replay should become a profile, not a
        niche configuration footnote.
\end{itemize}

\subsection{Non-Goals}

The design deliberately avoids three tempting but wrong shapes.
First, it does not reimplement Linux in another language. gVisor is
important prior art for backend structure and \KVM{} execution, but
\UML{} keeps Linux semantics by building Linux itself
\cite{gvisor-platforms}. Second, it is not a microVM competitor to
Firecracker \cite{firecracker}; the problem is kernel research and
fuzzing fidelity, not cloud cold-start packaging. Third, it is not a
unikernel-only library effort, although the library profile borrows
from LKL where that model is useful \cite{lkl}.

\subsection{Success Criteria}

The project succeeds when a developer can select a profile, build it
from the same tree, and get an artifact whose performance and
instrumentation properties match its intended use:

\begin{table}[h]
\centering
\caption{Target profile outcomes.}
\begin{tabularx}{\linewidth}{@{}p{0.21\linewidth}X@{}}
\toprule
Profile & Intended outcome \\
\midrule
\texttt{prod-fast} & \KVM{} backend, no hooks, near-native syscall path. \\
\texttt{prod-with-hooks} & Fast production binary with hooks compiled but off; tracing can be enabled after a crash without rebuilding. \\
\texttt{research} & Debug-friendly \UML{} with tracing, kprobes, sanitizers, and interactive control surfaces. \\
\texttt{fuzz} / \texttt{fuzz-deep} & KCOV, snapshots, sanitizer coverage, replay-oriented hooks, and syzkaller integration. \\
\texttt{sandbox} & Minimal debug surface and minimized host-side trust boundary. \\
\texttt{library} & LKL-style direct-call mode for harnesses that can sacrifice ring-split fidelity. \\
\texttt{time-travel} & Deterministic time and replay hooks as a first-class shipped configuration. \\
\bottomrule
\end{tabularx}
\end{table}

\subsection{Why Now}

The timing is favorable because recent \UML{} work removed old
prerequisite blockers: seccomp mode exists, SMP support exists, and
modern kernel instrumentation is mature enough to compose around
static keys and profile-specific Kconfig. The surrounding ecosystem
also supplies useful models: gVisor for backend/platform design,
crosvm for host process decomposition \cite{crosvm}, syzkaller for
fuzzing expectations \cite{syzkaller}, and upstream kernel policy for
AI-assisted development disclosure \cite{coding-assistants}.
